\chapter{Downtime Management}
\label{ch:downtime}
\index{downtime}

Downtime is the narrative space between adventures – days, weeks, or even months when characters rest, train, build assets, settle debts, and pursue personal projects. Unlike many games, \textit{Fate’s Edge} treats downtime as an active phase where players make meaningful choices and the GM introduces complications. This chapter gives you the tools to run downtime efficiently and dramatically.

\section{The Downtime Framework}
\label{sec:downtime-framework}

Each downtime phase follows a simple loop:

\begin{enumerate}
\item \textbf{Set the Duration.} Determine how much time passes (usually 1–4 weeks, but can be longer for epic campaigns).
\item \textbf{Player Declarations.} Each player states one or two major goals (training, research, carousing, etc.). See the table below for common actions.
\item \textbf{Roll + Resolve.} Call for appropriate rolls (Attribute + Skill) against DVs you set. Use the \textbf{Outcome Matrix} – a Miss or Partial can generate Story Beats or advance complication clocks.
\item \textbf{Upkeep \& Clocks.} Players pay upkeep for assets and followers, advance or retreat campaign clocks, and clear Obligation through off‑screen acts of service.
\item \textbf{Complications \& Twists.} Spend any Story Beats generated during downtime to introduce rumours, rival actions, or unexpected events that will drive the next session.
\item \textbf{Log the Results.} Use the downtime log template (see \ref{sec:downtime-log}) to record what happened for easy reference.
\end{enumerate}

\subsection{Common Downtime Rolls}
\label{subsec:common-rolls}

Use this table as a quick reference. The DV reflects typical effort; increase DV by +1 or +2 for rushed or dangerous projects.

\renewcommand{\arraystretch}{1.2}
\begin{longtable}{p{4cm}p{2.5cm}p{3cm}p{5cm}}
\toprule
\textbf{Goal} & \textbf{Skill Roll} & \textbf{Base DV} & \textbf{Effect on Success} \\
\midrule
\endfirsthead
\multicolumn{4}{c}{{\bfseries Continued from previous page}} \\
\toprule
\textbf{Goal} & \textbf{Skill Roll} & \textbf{Base DV} & \textbf{Effect on Success} \\
\midrule
\endhead
\bottomrule
\endfoot
Recover from Harm & Medicine + Wits & 3 & Remove 1 Harm (stacks with rest) \\
Train a Skill & relevant Attribute + Skill & 4 & Advance 1 segment on a Training Clock [4]. \\
Learn a New Talent & appropriate Attribute + Lore & 5 & Advance 1 segment on a Talent Acquisition Clock [6]. \\
Clear Obligation & Spirit + Lore (act of service) & 3 & Reduce Obligation by 1 segment (GM describes act). \\
Carouse for Information & Presence + Streetwise & 3 & Gain a rumour or a minor String. \\
Craft an Item & Wits + Craft & DV by item (3–6) & Create a mundane or minor magical item. \\
Research a Topic & Wits + Investigation & 4 & Gain one piece of true information (GM may be cryptic). \\
Manage an Asset & Presence + Command & 3 & Reduce or prevent condition degradation. \\
Establish a Contact & Presence + Sway & 4 & Gain a new Minor String. \\
Reduce a Debt Clock & Presence + Sway or service roll & 4 & Clear 1–2 segments (GM discretion). \\
Train a Follower & Spirit + Command & 4 & Improve Follower Cap by 1 (requires Training Clock). \\
\end{longtable}

\paragraph{Note on Training Clocks.} For long‑term improvement (skills, talents, follower advancement), start a Progress Clock [4–8]. Each successful downtime roll advances the clock by 1 segment. When full, the improvement is complete. Failure adds a complication instead of advancing.

\section{Rush Rules – Risk Clocks for Accelerated Development}
\label{sec:rush-rules}
\index{Rush rule}\index{training}\index{Risk Clock}

Players can attempt to \textbf{rush} a downtime activity – train a skill in half the normal days, craft an item overnight, or clear a debt in a single desperate scene. Rushing replaces the normal training time with a \textbf{Risk Clock} that fills as the player pushes too hard.

\subsection{Procedure}
\label{subsec:rush-procedure}
\begin{enumerate}
\item Player declares they are rushing (e.g., “I want to learn Acrobatics in 3 days instead of a week.”).
\item The GM creates a \textbf{Risk Clock} of 4–6 segments, named after the potential consequence (e.g., \textit{Backlash [4]}, \textit{Injury [6]}, \textit{Patron Ire [4]}).
\item For each downtime action related to the rush, the player rolls as normal. On a \textbf{Partial} or \textbf{Miss}, the clock advances by 1 (or more, at GM discretion). On a \textbf{Clean Success}, it does not advance.
\item \textbf{When the Risk Clock fills}, the rushed project fails or succeeds at a heavy cost. The GM chooses one:
  \begin{itemize}
    \item The project fails entirely, but the player gains 2 Boons to use on a later attempt.
    \item The project succeeds, but the character suffers a permanent Condition (e.g., \textit{Twisted Ankle}, \textit{Migraines}) that lasts until treated.
    \item The character gains the new ability, but immediately marks 1 Harm (overexertion) and gains a Debt Clock [4] to a terrestrial patron who “helped” them cheat time.
  \end{itemize}
\item If the clock never fills, the rush succeeds as normal with no extra cost.
\end{enumerate}

\begin{tcolorbox}[colback=red!5!white, colframe=red!75!black, title=Example: Rushing a Skill]
Kael wants to raise his Melee from 3 to 4. Normal training takes 4 days. He rushes to do it in 2 days. The GM starts a \textit{Pulled Muscle [6]} clock. Kael rolls three downtime actions (one per day). He gets a Partial and a Miss – the clock fills. The GM rules: “You gain the skill, but during the final drill you tear a tendon. Mark Harm 1 (Level 1) and you have the \textit{Limp} Condition until you receive magical healing or rest for a full week.”
\end{tcolorbox}

\section{Using Downtime to Advance Clocks}
\label{sec:downtime-advance-clocks}
\index{clocks!downtime}

Downtime is a perfect time for players to actively influence the campaign’s \textbf{Mandate}, \textbf{Crisis}, or \textbf{faction clocks}. Instead of just rolling for personal gain, they can declare they are working on a larger goal.

\subsection{Mandate and Crisis}
\label{subsec:mandate-crisis}
\begin{itemize}
\item \textbf{Increase Mandate (public favour):} Players can spend a downtime action to perform charity, solve a local problem, or build alliances. Roll Presence + Sway (DV 4). On success, increase Mandate by 1 (to a max of 6). On a Miss, increase Crisis by 1 (enemies notice their growing influence).
\item \textbf{Decrease Crisis (reduce threat):} Players can sabotage enemy plans, broker a truce, or expose a conspiracy. Roll appropriate skill (e.g., Stealth, Subterfuge, Lore) DV 5. On success, reduce Crisis by 1. On a Miss, Mandate also decreases by 1 (collateral damage).
\end{itemize}

\subsection{Faction and Debt Clocks}
\label{subsec:faction-debt-clocks}
\begin{itemize}
\item Players can take actions to tick \textbf{forward} or \textbf{backward} a visible faction clock. For example: “I spend the week bribing officials in the Dyers’ Guild.” Roll Presence + Sway (DV 4). On success, tick the \textit{Guild Influence} clock \textbf{back} by 1 (less threat). On a Miss, the clock ticks \textbf{forward}.
\item Accepting a Debt Clock from a terrestrial patron during downtime can be resolved with a service scene (roll to clear 1–2 segments) or by spending XP.
\end{itemize}

\section{Handling Off‑Screen Characters (Echo Characters)}
\label{sec:off-screen-characters}
\index{Echo character}\index{off-screen}

When a player has two or more characters (or a character steps off‑screen for multiple sessions), you need a lightweight way to keep them active in the fiction without splitting spotlight. Use the \textbf{Echo Character} system.

\subsection{Core Principles}
\begin{itemize}
\item Off‑screen characters are \textbf{not frozen}. They continue to act, gather information, and face consequences.
\item The player provides a brief \textbf{Downtime Log} (see \ref{sec:downtime-log}) for each off‑screen character.
\item You (the GM) spend 2–3 minutes per log, describing outcomes and possibly ticking a personal \textbf{Crisis Clock [6]} for that character.
\end{itemize}

\subsection{Switch Mechanics}
\label{subsec:switch-mechanics}
When a player wants to bring an off‑screen character back into play (or send a character off‑screen):
\begin{enumerate}
\item \textbf{Declare} – “I’m switching to my other character, the scholar, for the next session.”
\item \textbf{Status} – The current character becomes an Echo. They take one of three statuses:
  \begin{itemize}
    \item \textbf{Supporting} – Works toward active PC goals (gathers intel, manages assets). Minimal risk.
    \item \textbf{Advancing} – Pursues their own agenda (research, training). May uncover opportunities or trigger a complication.
    \item \textbf{In Peril} – Their absence has made them a target. The GM ticks a personal Crisis Clock each session until they return.
  \end{itemize}
\item \textbf{Re‑entry} – When returning, the character gains 2 Boons for a dramatic entrance. The player rolls Wits + Spirit (DV 3) to “re‑establish presence.” On a Miss, a complication follows them back.
\end{enumerate}

\subsection{Downtime Log Template}
\label{sec:downtime-log}
\index{downtime log}

Provide players with a simple log to fill out between sessions for each off‑screen (or even on‑screen) character. You can hand out a card or use this template in your session notes.

\begin{tcolorbox}[colback=gray!5!white, colframe=black, title=Downtime Log – \_\_\_\_\_\_\_\_\_ (Character Name)]
\begin{tabular}{p{0.9\textwidth}}
\textbf{Resource Action:} (How they maintain connections, pay upkeep, or keep assets running) \\
\rule{0pt}{4ex}\\
\textbf{Project Action:} (One goal they worked toward – training, research, crafting) \\
\rule{0pt}{4ex}\\
\textbf{Event Response:} (How they react to a rumour, threat, or off‑screen event you describe) \\
\rule{0pt}{4ex}\\
\textbf{GM Response:} (You fill this in after the session) \\
\rule{0pt}{4ex}
\end{tabular}
\end{tcolorbox}

\paragraph{Example Log (Player Filled):}
\begin{verbatim}
Resource Action: Paid upkeep for my safehouse (2 XP).
Project Action: Researched the Drowned Man (Wits+Lore, DV 4) -> rolled a Partial (1 SB to GM).
Event Response: My contact in the Scarlets told me Mira is moving cargo through the old dyeworks.
GM Response: The research reveals the Drowned Man was a dyer; the SB will spawn a rival who wants the same information. Mira’s cargo is actually a trapped spirit – you can intercept it next session.
\end{verbatim}

\subsection{Off‑Screen Crisis Clock}
Each off‑screen character tracks a personal \textbf{Crisis Clock [6]}. It ticks when:
\begin{itemize}
\item The player rolls a Miss on a downtime action.
\item The character’s status is \textit{In Peril}.
\item You (GM) spend a Story Beat to introduce trouble for that character.
\end{itemize}
When the clock fills, the character suffers a major off‑screen consequence: a rival strikes, an asset is destroyed, or they are captured. This creates immediate hooks for the active party.

\section{Additional Downtime Tools and Advice}
\label{sec:downtime-advice}

\subsection{Asset \& Follower Upkeep}
Remind players to pay upkeep (Efficient or Intensive) for each asset and follower. Use the templates from Chapter 10 (Assets and Followers). If they forget, degrade the asset to \textit{Neglected} or a follower to \textit{Wary}. You can also use the “Rush” rules to restore a degraded asset in half the normal downtime.

\subsection{Resolving Complications}
Downtime is an excellent moment to clear lingering Complications (from character creation or story). Players can spend a downtime action to “resolve a complication” – roll appropriate skill (DV 3–5). On success, the complication is cleared; on a Miss, it worsens (GM ticks a related clock).

\subsection{Using Story Beats from Downtime Rolls}
Any 1s rolled during downtime generate Story Beats. Do not waste them. Use the SB menu:
\begin{itemize}
\item \textbf{1 SB:} A minor rumour begins about the character.
\item \textbf{2 SB:} A rival or faction becomes aware of their activities.
\item \textbf{3 SB:} A complication arises – a theft, a broken item, a social faux pas.
\item \textbf{4+ SB:} A major twist – a Patron demands an audience, an enemy strikes an asset.
\end{itemize}
These complications will ignite the next session’s opening scene.

\subsection{Downtime as Collaborative Worldbuilding}
Encourage players to ask questions during downtime: “I try to find a fence who deals in cursed coins. Is there one in this city?” Use their rolls to establish new NPCs and locations. A Clean Success means the fence exists and is reliable; a Partial means they exist but have a price (Debt Clock); a Miss means they exist but are working for an enemy.

\subsection{Sample Downtime Sequence (GM Script)}
\begin{quote}
“Two weeks pass. You each have two downtime actions. I need: a roll for skill training, a roll for asset upkeep, and a roll for whatever personal project you declared. You can also use a downtime action to interact with the campaign clocks – tell me how you want to affect the Mandate or the Cult’s Progress clock. Describe your goal, I’ll set a DV, and we’ll roll. Any 1s will give me Story Beats to complicate your return to action. Ready?”
\end{quote}

\section{Faction Turn Cheat Sheet}
\label{sec:faction-turn}
\index{factions!turn}\index{clocks!faction}

Between sessions (or at the start of a session), you can run a quick **faction turn** to keep the world alive. A faction turn takes 5–10 minutes and answers: *“What have the major powers done while the party was busy?”*

\subsection{When to Run a Faction Turn}
\begin{itemize}
  \item At the end of a session, after players declare their next destination.
  \item At the beginning of a session, as a “news from the world” recap.
  \item During a long downtime (several weeks or months).
\end{itemize}

\subsection{Procedure – 5 Minutes}
\begin{enumerate}
  \item \textbf{List active factions.} Keep no more than 3–5 that directly interact with the party.
  \item \textbf{For each faction, roll 1d6.} Use the table below to determine their progress.
  \item \textbf{Advance or retreat their Agenda Clock.} Each faction has a visible clock (6–10 segments) representing their goal.
  \item \textbf{Note one concrete change.} Write one sentence per faction (e.g., “The Scarlets seized the river docks.”).
  \item \textbf{Read the changes aloud} at the start of the next session.
\end{enumerate}

\begin{tcolorbox}[colback=blue!5!white, colframe=blue!75!black, title=Faction Progress Roll (d6)]
\begin{center}
\begin{tabular}{cll}
\toprule
\textbf{Roll} & \textbf{Progress} & \textbf{Agenda Clock Change} \\
\midrule
1–2 & Setback & –1 segment (min 0) \\
3–4 & Stalled & No change \\
5–6 & Advance & +1 segment \\
\bottomrule
\end{tabular}
\end{center}
\end{tcolorbox}

\subsection{Optional: Faction Conflict Roll}
If two factions are directly opposed (e.g., the Scarlets and the Grey Circle), roll opposed d6. Higher roll wins; lower roll loses 1 segment. Ties mean a stalemate – both lose 1 segment.

\subsection{Example Faction Turn}
\begin{quote}
\textbf{Factions:}
\begin{itemize}
  \item Mira’s Scarlets – Agenda Clock: \textit{Control the Dye Trade} [8] (currently 3/8)
  \item Church of the Everflame – Agenda Clock: \textit{Suppress Heresy} [6] (currently 2/6)
  \item The Drowned Man (not a faction, but a wildcard) – no clock
\end{itemize}

\noindent \textbf{Rolls:}
\begin{itemize}
  \item Scarlets: 5 → Advance → now 4/8.
  \item Everflame: 2 → Setback → now 1/6.
\end{itemize}

\noindent \textbf{One‑sentence changes:}
\begin{itemize}
  \item *Scarlets:* “Mira’s agents have bought the loyalty of the harbour master.”
  \item *Everflame:* “An inquisitor was found murdered; the Church retreats to its cathedral.”
  \item *Drowned Man:* (GM fiat) “The Drowned Man dragged two Scarlets into the canal. Rumours spread.”
\end{itemize}

\noindent \textbf{Next session opener:} “You wake to news that the harbour master has doubled tariffs – and that a body was pulled from the water wearing a Scarlets’ token.”
\end{quote}

\subsection{Using Faction Clocks in Play}
\begin{itemize}
  \item \textbf{When a clock fills}, the faction achieves its goal. The party must react (or suffer the consequences).
  \item Players can \textbf{affect faction clocks} during adventures: a successful sabotage mission might tick a clock backwards; a public failure might tick it forward.
  \item Keep faction clocks \textbf{visible} (on an index card or whiteboard). Players should see the pressure building.
\end{itemize}

\subsection{Lightweight Faction Standing (Optional)}
For factions that are not pursuing a single agenda but merely have an attitude toward the party, track a simple \textbf{-3 to +3 scale}:
\begin{itemize}
  \item \textbf{-3:} Enemy – actively works against the party.
  \item \textbf{0:} Neutral – indifferent.
  \item \textbf{+3:} Ally – will sacrifice for the party.
\end{itemize}
Shift standing by ±1 after major interactions (saving a faction’s leader, betraying them, etc.). No dice needed – just fiction.

\subsection{GM Quick Cues}
\begin{itemize}
  \item \textbf{Don’t overdo it.} One faction turn per 1–2 sessions is plenty.
  \item \textbf{Tie changes to player actions.} If the party ignored a threat, tick that faction’s clock for free.
  \item \textbf{Use the turn to generate rumours.} “The Grey Circle is digging a new well – some say they found bones.”
\end{itemize}

\section{Downtime at a Glance – GM Cheat Sheet}
\label{sec:downtime-cheat-sheet}

\begin{tcolorbox}[colback=blue!5!white, colframe=blue!75!black, title=Downtime Quick Reference]
\begin{itemize}
\item \textbf{Set duration} (1–4 weeks typical).
\item \textbf{Players declare} 1–2 goals (see table of common rolls).
\item \textbf{Roll} Attribute+Skill vs DV (3–6).
\item \textbf{Apply outcome}: Clean Success = progress; Partial = progress + SB; Miss = no progress + SB + maybe Risk Clock tick.
\item \textbf{Spend SB} to introduce rumours, rival actions, or complications.
\item \textbf{Advance or retreat} Mandate, Crisis, faction clocks based on player choices.
\item \textbf{Handle off‑screen characters} via Downtime Logs and Crisis Clocks.
\item \textbf{Rush Rule}: Use Risk Clock [4–6] for accelerated training/crafting.
\item \textbf{Asset/Follower Upkeep}: Players pay XP or use downtime actions; degrade if forgotten.
\item \textbf{Clear Complications} with a successful roll.
\end{itemize}
\end{tcolorbox}

\noindent With these tools, downtime becomes a productive, tension‑filled phase of the game rather than a mere pause. The world keeps moving, and the players’ choices during rest directly shape the next adventure.

